% Компактное описание проекта (3--4 стр.): ответы на вопросы чек-листа

\chapter{{\fontsize{14pt}{21pt}\selectfont\MakeUppercase{Описание проекта}}}
\label{ch:body}

\section{Актуальность, цель и задачи}

\textbf{Актуальность.} Значительная доля пользователей в России использует геосервисы для поиска офлайн-бизнеса; лидирующие платформы (Яндекс.Карты, 2ГИС) не предоставляют публичного API. Существующие решения (CRM, SMM, ORM) не закрывают задачу консистентности данных на разнородных каналах и избыточны для МСБ.

\textbf{Цель работы:} разработать программный комплекс на основе мультиагентной архитектуры с гибридной моделью интеграции (API~+ RPA) для управления цифровым присутствием субъектов малого и среднего бизнеса.

\textbf{Задачи} (в хронологическом порядке):
\begin{enumerate}
  \item Проанализировать предметную область и существующие решения, доступность API целевых платформ.
  \item Спроектировать классификацию платформ по типу интеграции (API / RPA) и платформонезависимую архитектуру (оркестратор, агенты).
  \item Разработать API-агенты (VK, Telegram) и RPA-агент (Яндекс.Бизнес, Playwright) на Go.
  \item Разработать веб-интерфейс на Next.js для управления настройками и мониторинга.
  \item Протестировать API- и RPA-агенты на тестовых экземплярах платформ, оценить характеристики.
\end{enumerate}

\section{Проблема и целевая аудитория}

\textbf{Какую проблему решает проект?} МСБ вынужден вручную поддерживать актуальную и согласованную информацию на множестве платформ (Яндекс.Карты, 2ГИС, VK, Telegram). Это даёт рассогласование данных, потерю клиентов и репутационные риски. Часть платформ не имеет публичного API (Яндекс.Бизнес, 2ГИС), что исключает простую автоматизацию.

\textbf{У кого возникает проблема?} Субъекты малого и среднего бизнеса с офлайн-точкой присутствия. Типовой пример~--- кофейня: малое предприятие, для которого отдельный сотрудник на цифровое присутствие нецелесообразен; критичны часы работы, адрес, отзывы.

\section{Как проблема решается в настоящее время}

Ручное управление на каждой платформе: отдельная авторизация, последовательное обновление данных по площадкам, поочерёдная проверка отзывов и сообщений. Отсутствует эталонный источник данных и централизованное журналирование. Альтернативы~--- комплексные CRM/SMM (Bitrix24, Hootsuite и~др.), но они не решают консистентность данных и часто недоступны по цене для малого бизнеса.

\section{Аналоги, конкуренты и критерии сравнения}

\textbf{Аналоги:} CRM и омниканальность (Bitrix24, SendPulse, Kommo), SMM-планирование (Hootsuite, Buffer, SMMplanner), мониторинг и ORM (YouScan). Ни одно решение не обеспечивает эталонных данных и контроля расхождений между каналами, единого протокола для API и RPA, интеграции с Яндекс.Бизнес и 2ГИС.

\textbf{Критерии сравнения} (метод взвешенных баллов 0--3): функциональный охват, мультиканальность, автоматизация, управление консистентностью данных (ключевой), трассируемость, расширяемость. OneVoice проектируется с фокусом на консистентности и гибридной интеграции (проектный итог выше аналогов).

\section{Предметная область (исследовательская часть)}

Цифровое присутствие~--- совокупность данных о компании в цифровых каналах (карточки на картах, соцсети, мессенджеры, отзывы). Ключевые контуры: консистентность справочных данных и репутация (отзывы, ответы). Исследовательская часть: классификация платформ по типу интеграции (API~--- VK, Telegram; RPA~--- Яндекс.Бизнес, 2ГИС); обоснование мультиагентного подхода и протокола A2A.

\section{Практическая часть: технологии и инструменты}

Языки и фреймворки: \textbf{Go} (сервер, агенты), \textbf{Next.js} + React + Tailwind + shadcn/ui (клиент), \textbf{PostgreSQL} и \textbf{MongoDB} (данные), \textbf{NATS} (очереди), \textbf{Playwright} (RPA). Протокол взаимодействия агентов~--- A2A поверх NATS. Агенты MVP: VK (API), Telegram (API), Яндекс.Бизнес (RPA).

\section{Данные: источники и персональные данные}

Данные: профили бизнеса, расписания, интеграции, диалоги, задачи, отзывы, публикации. Источники: ввод пользователя, API (VK, Telegram), RPA (Яндекс.Бизнес). Персональные данные обрабатываются в соответствии с 152-ФЗ: пароли (bcrypt), токены платформ (AES-256-GCM), разграничение доступа (RBAC), аудит, без чувствительных данных в логах.

\section{Итог}

Проект OneVoice направлен на автоматизацию управления цифровым присутствием МСБ при гетерогенности платформ (API + RPA). Цель и задачи сформулированы; проблема, ЦА, аналоги, предметная область, технологии и данные описаны в соответствии с чек-листом описания проекта.
